\section{Review Exercises}

\subsection{Elementary computational and construction exercises}

\begin{enumerate}
\item Let $A=(1,2)$ and $B=(5,-1)$. Compute the displacement vector $\overrightarrow{AB}$ and draw the same vector with tail at the origin and with tail at $P=(-2,3)$.

\item For
\[
u=(3,-2),\qquad v=(-1,4),
\]
compute $u+v$, $u-v$, $2u$, and $-3v$. Draw $u$, $v$, and $u+v$ using the parallelogram rule.

\item Compute the linear combination
\[
2u-3v
\]
for $u=(1,2)$ and $v=(2,-1)$. Verify the result by carrying out the corresponding geometric sequence of displacements.

\item Compute the length of each vector:
\[
u=(3,4),\qquad v=(-5,12),\qquad w=(2,-2).
\]
Then find the corresponding unit vector in each direction.

\item For
\[
u=(2,1),\qquad v=(1,-2),\qquad w=(4,2),
\]
compute all three dot products $u\cdot v$, $u\cdot w$, and $v\cdot w$. Decide which pairs are perpendicular and which are parallel.

\item Find the angle between
\[
u=(1,0),\qquad v=(1,1).
\]
Then find the angle between $v=(1,1)$ and $w=(-1,1)$.

\item Compute the projection of
\[
u=(4,3)
\]
onto $v=(1,0)$ and onto $w=(1,1)$. For each case, write the orthogonal decomposition of $u$ into a component parallel to the chosen vector and a perpendicular remainder.

\item Let
\[
b_1=(1,1),\qquad b_2=(1,-1),\qquad u=(4,2).
\]
Find the coordinates $(\alpha,\beta)$ such that
\[
u=\alpha b_1+\beta b_2.
\]
Check the answer by reconstruction.

\item The same vector has basis coordinates $[u]_B=(2,-1)$ in the basis
\[
B=(b_1,b_2),\qquad b_1=(1,2),\quad b_2=(2,1).
\]
Find its Cartesian coordinates. Then solve for its coordinates in the standard basis directly.

\item Convert between Cartesian and polar coordinates:
\begin{enumerate}
\item find polar coordinates of $(\sqrt3,1)$ and $(-1,\sqrt3)$;
\item find Cartesian coordinates of $(r,\theta)=(4,\pi/6)$ and $(3,\pi)$.
\end{enumerate}
\end{enumerate}

\subsection{Conceptual and reasoning exercises}

\begin{enumerate}
\item Explain the difference between a point and a vector. Why can the same vector be drawn from different starting points while a point cannot be translated without becoming a different point?

\item Prove from the geometric meaning of successive displacements that vector addition is associative. Do not begin from the coordinate formula.

\item A student claims that $u+v$ must always be longer than both $u$ and $v$. Give a counterexample and explain geometrically what causes the failure.

\item Explain why scalar multiplication by a negative number changes both length and direction. What remains unchanged when a vector is multiplied by any nonzero scalar?

\item Show that two nonzero vectors are parallel exactly when one is a scalar multiple of the other. Explain how the coordinate criterion expresses a geometric fact.

\item Use the cosine theorem to explain why the dot-product formula
\[
u\cdot v=\|u\|\|v\|\cos\theta
\]
contains information about both lengths and angle.

\item Explain why $u\cdot v=0$ characterises perpendicularity only when both vectors are interpreted in Euclidean geometry. What geometric role is played by the common coordinate scale?

\item A basis changes but the geometric vector does not. Explain precisely what changes and what stays fixed. Why are basis coordinates properties of the pair ``vector plus basis'' rather than of the vector alone?

\item Explain why two nonparallel vectors form a basis of the plane. What fails if the two proposed basis vectors are parallel?

\item Compare Cartesian and polar coordinates as two languages for the same point or displacement. Give one kind of problem for which each system reveals the geometry more naturally.
\end{enumerate}
